\documentclass[letterpaper,11pt]{article}

% --- PAQUETES DE IDIOMA
\usepackage[spanish,mexico]{babel}
\usepackage[utf8]{inputenc}
\usepackage{csquotes}
\usepackage{iftex}

\usepackage{pdfpages}

% --- PAQUETES MATEMÁTICOS --
\usepackage{amsmath}
\usepackage{cancel}
\usepackage{mathrsfs}

% --- PAQUETES GRÁFICOS Y DE DISEÑO ---
\usepackage{graphicx}
\usepackage{float}
\usepackage{placeins}
\usepackage{subcaption}
\usepackage[table]{xcolor}
\usepackage{makecell}
\usepackage{array}
\usepackage{tabularx}
\usepackage{longtable}
\newcolumntype{Y}{>{\raggedright\arraybackslash}X}
\usepackage{tikz}
\usetikzlibrary{shapes.geometric, shapes.multipart, positioning, arrows.meta, calc, shadows, decorations.pathreplacing, fit, babel}
\usepackage{forest}
\usepackage{pgfgantt}

% --- LISTAS ---
\usepackage{enumitem}

% --- PAQUETES DE CÓDIGO ---
\usepackage{listings}
% Configuración del estilo de listings
\definecolor{codebg}{rgb}{0.95,0.95,0.95}   
\definecolor{keyword}{rgb}{0.0,0.2,0.6}     
\definecolor{comment}{rgb}{0.25,0.5,0.35}   
\definecolor{string}{rgb}{0.6,0.0,0.0}      
\definecolor{number}{rgb}{0.5,0.0,0.5}      
\definecolor{codegray}{rgb}{0.5,0.5,0.5}
\definecolor{codepurple}{rgb}{0.58,0,0.82}
\definecolor{backcolour}{rgb}{0.95,0.95,0.95}

\lstset{
  backgroundcolor=\color{backcolour},
  basicstyle=\ttfamily\small,
  keywordstyle=\color{blue}\bfseries,
  stringstyle=\color{codepurple},
  commentstyle=\color{codegray}\itshape,
  numbers=left,
  numberstyle=\tiny\color{codegray},
  frame=single,
  rulecolor=\color{black},
  breaklines=true,
  tabsize=4,
  showstringspaces=false,
  literate={á}{{\'a}}1
           {é}{{\'e}}1
           {í}{{\'i}}1
           {ó}{{\'o}}1
           {ú}{{\'u}}1
           {Á}{{\'A}}1
           {É}{{\'E}}1
           {Í}{{\'I}}1
           {Ó}{{\'O}}1
           {Ú}{{\'U}}1
           {ñ}{{\~n}}1
           {Ñ}{{\~N}}1
           {¿}{{\textquestiondown}}1
           {¡}{{\textexclamdown}}1
}

% Estilo para código en C
\lstdefinestyle{cstyle}{
  language=C,
  backgroundcolor=\color{backcolour},
  basicstyle=\ttfamily\small,
  keywordstyle=\color{blue}\bfseries,
    keywordstyle=[2]\color{yellow!70!orange}\bfseries,
    morekeywords=[2]{MPI_Init,MPI_Comm_size,MPI_Comm_rank,MPI_Send,MPI_Recv,MPI_Scatter,MPI_Gather,MPI_Bcast,MPI_Reduce,MPI_Finalize},
    emph={printf},
    emphstyle=\color{blue}\bfseries,
  stringstyle=\color{codepurple},
  commentstyle=\color{codegray}\itshape,
  numbers=left,
  numberstyle=\tiny\color{codegray},
  frame=single,
  rulecolor=\color{black},
  breaklines=true,
  tabsize=4,
  showstringspaces=false
}

% Definición del lenguaje Python para listings
\lstdefinelanguage[LAB]{Python}{
    morekeywords={False, None, True, and, as, assert, async, await, break,
                                class, continue, def, del, elif, else, except, finally,
                                for, from, global, if, import, in, is, lambda, nonlocal,
                                not, or, pass, raise, return, try, while, with, yield,
                                machine, Pin, PWM, ADC, Timer, sleep_ms, sleep_us},
    morecomment=[l]{\#},
    morestring=[b]",
    morestring=[b]',
    sensitive=true
}

% Estilo para salidas de consola/terminal
\lstdefinestyle{output}{
  backgroundcolor=\color{black!5},
  basicstyle=\ttfamily\footnotesize,
  numbers=none,
  frame=single,
  rulecolor=\color{black!30},
  breaklines=true,
  tabsize=4,
  showstringspaces=false,
  xleftmargin=0.5cm,
  xrightmargin=0.5cm
}

% --- ESTILOS PARA DIAGRAMAS DE FLUJO (TIKZ) ---
\tikzset{
  startstop/.style={rectangle, rounded corners, minimum width=3cm, minimum height=1cm,text centered, draw=black, fill=red!30},
  process/.style={rectangle, minimum width=3cm, minimum height=1cm, text centered, draw=black, fill=orange!30},
  decision/.style={diamond, aspect=2, minimum width=3cm, minimum height=1cm, text centered, draw=black, fill=green!30},
  io/.style={trapezium, trapezium left angle=70, trapezium right angle=110, minimum width=3cm, minimum height=1cm, text centered, draw=black, fill=blue!30},
  arrow/.style={thick,->,>=stealth}
}


\definecolor{azulresultado}{HTML}{0072C6}
\definecolor{rojoacarreo}{HTML}{FF0000}

% --- RUTAS DE IMÁGENES ---
\graphicspath{{img/}{portada_img/}}

\definecolor{armblue}{RGB}{0, 113, 188}

% --- ÍNDICE CON PUNTOS ---
\usepackage{tocloft}
\renewcommand{\cftsecleader}{\cftdotfill{\cftdotsep}}

% --- CONFIGURACIÓN DE PÁGINA (GEOMETRÍA) ---
\usepackage[
    left=25mm, 
    right=25mm,
    top=35mm,
    bottom=30mm,
    headheight=77pt, 
    papersize={21.59cm,27.94cm}
]{geometry}

% --- FORMATO DE TEXTO ---
\usepackage{setspace}
\onehalfspacing
\usepackage{parskip} 
\usepackage{lipsum}  

% --- VARIABLES GLOBALES ---
\newcommand{\materia}{}

% --- ENCABEZADOS Y PIES DE PÁGINA ---
\usepackage{fancyhdr}
\pagestyle{fancy}
\fancyhf{} 
\fancyhead[L]{\includegraphics[height=1.5cm]{portada_img/der.png}}
\fancyhead[R]{\includegraphics[height=1.5cm]{portada_img/izq.png}}
\fancyfoot[C]{\thepage} 
\renewcommand{\headrulewidth}{0.4pt} 
\renewcommand{\footrulewidth}{0pt}  

\fancypagestyle{indice}{
  \fancyhf{}
  \fancyhead[L]{\includegraphics[height=1.5cm]{portada_img/der.png}}
  \fancyhead[R]{\includegraphics[height=1.5cm]{portada_img/izq.png}}
  \renewcommand{\headrulewidth}{0.4pt}
  \renewcommand{\footrulewidth}{0pt}
}
\fancypagestyle{plain}{
  \fancyhf{}
  \fancyhead[L]{\includegraphics[height=1.5cm]{portada_img/der.png}}
  \fancyhead[R]{\includegraphics[height=1.5cm]{portada_img/izq.png}}
  \fancyfoot[C]{\thepage}
  \renewcommand{\headrulewidth}{0.4pt}
  \renewcommand{\footrulewidth}{0pt}
}

% --- BIBLIOGRAFÍA ---
\usepackage[backend=biber, style=apa, sortcites, url=true]{biblatex}
\addbibresource{referencias.bib}

% --- HIPERVÍNCULOS ---
\usepackage{hyperref}
\hypersetup{
    colorlinks=true,
    linkcolor=blue,
    citecolor=blue,
    filecolor=magenta,
    urlcolor=cyan
}

\begin{document}

\input{portada.tex}

\newpage
{%
  \makeatletter
  \let\ps@plain\ps@indice
  \makeatother
  \pagestyle{indice}
  \hypersetup{linkcolor=black}\tableofcontents
}


\newpage

\pagestyle{fancy}
\setcounter{page}{1}
\singlespacing

% ===============================================================
\section{Objetivo}
\subsection{Objetivo}

El presente proyecto tiene como objetiv principal, desarrollar una web, 
la cual permita a los usuarios comprar y vender cartas de Pókemon
y videojuegos fisicos de diferentes consolas, con esto 
haciendo un enfoque en marcas como Playstation, Xbox y Nintendo, mantiendo un modelo 
de negocio de compraventa entre usuarios, precios actulizados periodicamente mediante APIs externas,
trasacciones protegidas (integrando un escrow simulado), asi como la integracion de una logistica visual, la aplicacion estara diseñada 
para convertirse en un aliado estrategico para los usuarios que compran cartas pokemon, asi como a los usuarios que apoyen al formato fisico en los videojuegos.


\textbf{Objetivos de negocio}
\begin{enumerate}[label=\textbf{ON\arabic*.}, leftmargin=*]
    \item Dar una solucion al comercio diversificado de cartas 
    de pokemon, esto enfocado en México, con la finalidad de crear un espacio donde los coleccionistas puedan intercambiar sus cartas de manera segura y eficiente.
    \item Definir el modelo de ingresos (comisión por venta) y demostrar su cálculo automático dentro de la simulacion de compra.
    \item Establecer un punto de compra y venta para los videojuegos en formato fisico, esto para 
    los usuarios que esten interesados en adquirir videojegos de forma mas barata y para los usuarios que quieran vender sus videojuegos usados.
\end{enumerate}

\textbf{Objetivos de desarrollo web}
\begin{enumerate}[label=\textbf{OD\arabic*.}, leftmargin=*]
    \item Implementar una arquitectura cliente-servidor con consumo de APIs externas y actualización automatica de precios.
    \item Desarrollo de un Control de Acceso Basado en Roles para compradores, vendedores, recolectores y administradores.
    \item Integraciaso de un metodo de pago, esto por medio de el SDK de Mercado Pago.
\end{enumerate}

\newpage

% ===============================================================
\section{Definición del Producto o Servicio}
\subsection{Descripción del producto}


El producto es una plataforma web 
donde el objetivo es el comercio electrónico 
donde se va a mantener un enfoque en 
la compraventa de cartas Pokémon TCG (Esto teniendo en cuenta el numero 
de expanciones que se llevan con año) y videojuegos 
físicos (PS3, PS4, PS5, PSP, Xbox (original, 360 y X/S) y 
Nintendo (Nintendo Switch)), 
esto bajo la modalidad de comprar bajo vendedores particulares,
ahora bien, a diferencia de los mercados generalistas, la plataforma 
trata el coleccionismo como un mercado de valores, esto por medio 
de cada artículo muestra la evolución de su precio (esto se puede saber por medio de el mercado de valores de
PSA que existe actualmente), las condiciones de venta 
seran catalogadas en las cartas gradeadas, los datos de su grado y certificado, y las 
cartas sin gradear, donde son conocidas como (downGrade). 

Las transacciones se realizan con un esquema de compra protegida simulado 
y el comprador seguira el progreso del envio mismo.

La web ser responsiva y accesible desde el navegador, 
sin alguna instalación previa.

\subsection{Usuarios y roles}
\begin{table}[H]
\centering
\begin{tabularx}{\textwidth}{|l|Y|}
\hline
\textbf{Rol} & \textbf{Lo que hace} \\ \hline
Comprador & Encontrar artículos y decidir con datos: historial de precio, datos del grado PSA y compra protegida. \\ \hline
Vendedor (particular) & Publicar (Agregar descripcion, precio, generacion a la que pertenece y fotografias del producto por todos los lados posibles)y gestionar su inventario y consultar sus ventas. \\ \hline
Recolector & Actualizar los estados logísticos del pedido (rol operativo de la plataforma). \\ \hline
Administrador & Moderar publicaciones y usuarios, resolver disputas y supervisar métricas globales. \\ \hline
\end{tabularx}
\caption{Roles de la plataforma}
\end{table}

\newpage

\subsection{Requerimientos funcionales}

\begin{table}[H]
\centering
\begin{tabularx}{\textwidth}{|l|Y|}
\hline
\textbf{ID} & \textbf{Funcionalidad} \\ \hline
RF01 & Registro, inicio de sesión y asignación de roles. \\ \hline
RF02 & Catálogo de cartas (secciones sin gradear y gradeadas) y de videojuegos, con búsqueda y filtros. \\ \hline
RF03 & PPublicacion de los articulos con fotos de todos los lados posibles. \\ \hline
RF04 & Detalle con gráfica histórica de precio esto de las cartas y, en las gradeadas, grado, número de certificado. \\ \hline
RF05 & Compra con pago en sandbox de Mercado Pago. \\ \hline
RF06 & Panel del Recolector para actualizar estados del envío. \\ \hline
RF07 & Comparativa de precios con Mercado Libre. \\ \hline
RF08 & Vista de la logística del envío. \\ \hline
RF9  & Dashboard del vendedor y panel de administración. \\ \hline
RF10 & Wishlist con alertas de cambio de precio. \\ \hline
\end{tabularx}
\caption{Requerimientos funcionales}
\end{table}

\subsection{Estructura del catálogo: cartas sin gradear y gradeadas}

El catálogo de cartas se organiza en dos secciones. 

\begin{table}[H]
\centering
\small
\begin{tabularx}{\textwidth}{|l|Y|Y|Y|}
\hline
\textbf{Sección} & \textbf{Contenido} & \textbf{Condición} & \textbf{Precio de referencia} \\ \hline
Cartas sin gradear & Cartas sin certificar, descritas por el vendedor con fotografías. & Near Mint, Lightly Played, Moderately Played, Heavily Played, Damaged. & Valor de mercado de PokemonPriceTracker. \\ \hline
Cartas gradeadas & Cartas certificadas por PSA (de forma exclusiva solo se admitira de esta plataforma) 
y selladas en su acrilico correspondiente. & Grado numérico de la graduadora (en PSA, del 1 al 10). & Valor de mercado de PokemonPriceTracker. \\ \hline
Videojuegos & Videojuegos físicos por plataforma. & Nuevo sellado, completo en caja. & Precio fijado por el vendedor y comparativa con Mercado Libre. \\ \hline
\end{tabularx}
\caption{Secciones del catálogo}
\end{table}

\newpage
La condición de las cartas sin gradear se describe con cinco niveles:

\begin{enumerate}
    \item Near Mint (prácticamente sin uso, esquinas y bordes intactos)
    \item Lightly Played (desgaste mínimo, con pequeñas imperfecciones apenas visibles)
    \item Moderately Played (desgaste visible pero la carta sigue siendo legible y jugable)
    \item Heavily Played (desgaste notorio, con dobleces, rayones o manchas)
    \item Damaged (daño estructural importante, como rasgaduras o dobleces pronunciados)
\end{enumerate}

\textbf{Sección de PSA:} 

Dentro de las cartas gradeadas se tiene la sección 
para cartas PSA, 
porque PSA es la graduadora dominante en TCG, esto en el mercado se ve muy reflejado por que 
muchas personas es la casa graduadora que mas se utiliza.

Cada ficha PSA incluye el grado, el número de certificado y de su etiqueta, la gráfica de el historial de precios para ese grado.
\subsection{Fuentes de datos y servicios}
\begin{itemize}
    \item \textbf{PokemonPriceTracker}: valores de mercado de cartas, tanto sin gradear como gradeadas (incluidas las de PSA).
    \item \textbf{Pokémon TCG API}: catálogo de cartas.
    \item \textbf{Mercado Libre}: precios de referencia para la comparativa; se verificará si el endpoint de búsqueda requiere autenticación.
    \item \textbf{Videojuegos:} la fuente de precios sera PriceCharting. 
\end{itemize}

\subsection{Diferenciadores}
El proyecto cuenta con cuatro diferenciadores principales. En primer lugar, ofrece transparencia de precios al mostrar un historial construido mediante snapshots propios y una sección exclusiva para cartas gradeadas por PSA que incluye su respectivo enlace de verificación.

Además, garantiza una compra protegida C2C, ya que el sistema retiene el pago de forma segura hasta que el comprador confirma la recepción del artículo o hasta que vence el plazo de liberación de fondos.

En cuanto a la experiencia del usuario, la plataforma integra una logística visual que sustituye el tradicional estado estático por una animación del envío directamente sobre un mapa. Finalmente, todo el proyecto destaca por su bajo costo de infraestructura, debido a que el sistema se ejecuta íntegramente en contenedores Docker, eliminando por completo la dependencia de proveedores de nube de pago.

\section{Contexto del Negocio / Mercado}



\newpage
\section{Estrategia a seguir para el desarrollo}

\subsection{Estructura del trabajo (WBS)}
\begin{figure}[H]
\centering
\begin{forest}
for tree={font=\footnotesize, s sep=3pt, inner sep=1.5pt, grow'=0, child anchor=west, parent anchor=south, anchor=west, calign=first,
  edge path={\noexpand\path [draw, \forestoption{edge}] (!u.south west)+(7.5pt,0) |- node[fill,inner sep=1.25pt] {} (.child anchor)\forestoption{edge label};},
  before typesetting nodes={if n=1{insert before={[,phantom]}}{}},
  fit=band, before computing xy={l=15pt}}
[{Plataforma web TCG y Retro Gaming}
  [{1. Inicio y configuración}
    [{1.1 Conformación del equipo, roles, tablero y repositorio \textit{(VERSIONCTRL-1.0)}}]
    [{1.2 Especificación de requisitos (ERS) e historias de usuario \textit{(ERS-1.0)}}]
    [{1.3 Entornos Docker, CI y credenciales de APIs \textit{(DOCKER-1.0, APIEXT-1.0)}}]
    [{1.4 Diseño del esquema de BD (ER) y migraciones \textit{(DBMARIA-1.0)}}]
    [{1.5 Estudio de mercado y modelo de negocio (Canvas) \textit{(MODNEG-1.0)}}]
  ]
  [{2. Backend y APIs}
    [{2.1 Autenticación y RBAC}]
    [{2.2 Catálogo (sin gradear, gradeadas y PSA) y publicaciones \textit{(API-1.0 (parcial))}}]
    [{2.3 Adaptadores de APIs externas y scheduler de precios \textit{(PRECIOS-1.0)}}]
    [{2.4 Contrato OpenAPI y datos de ejemplo \textit{(OPENAPI-1.0)}}]
  ]
  [{3. Frontend}
    [{3.1 Prototipos UI y SPA base (sistema de diseño) \textit{(PROTO-1.0)}}]
    [{3.2 Catálogo por secciones, búsqueda y detalle con gráfica de precios}]
    [{3.3 Publicación de artículos (4 fotos) y paneles por rol}]
  ]
  [{4. Pagos, escrow y mapas}
    [{4.1 Mercado Pago (sandbox) y webhooks}]
    [{4.2 Máquina de estados del escrow y disputas \textit{(API-1.0)}}]
    [{4.3 Geocodificación Mapbox y mapa animado (Leaflet)}]
    [{4.4 Paneles de Recolector y Admin \textit{(WEB-1.0)}}]
  ]
  [{5. Integración y pruebas}
    [{5.1 Pruebas unitarias y de integración \textit{(CP-1.0)}}]
    [{5.2 Pruebas de seguridad y de carga \textit{(SEG-1.0)}}]
    [{5.3 Corrección de errores y validación con usuarios}]
  ]
  [{6. Despliegue y cierre}
    [{6.1 Empaquetado en Docker y respaldos \textit{(DEPLOY-1.0)}}]
    [{6.2 Documentación técnica y manual de usuario \textit{(MANUAL-1.0)}}]
    [{6.3 Cierre: demostración y entrega final \textit{(ENTREGA-1.0)}}]
  ]
]
\end{forest}
\caption{WBS}
\end{figure}

Se emplea la estimación por tres valores (PERT), porque el proyecto integra varios servicios externos y tecnologías nuevas para el equipo, con factores que pueden retrasar las actividades (cuotas de las APIs, limitaciones del \textit{sandbox} y del hardware). Para cada actividad se estiman la duración optimista (O), la más probable (M) y la pesimista (P); la duración esperada es $E=(O+4M+P)/6$ y la varianza es $\sigma^2=((P-O)/6)^2$. La duración programada es el valor de $E$ redondeado al día entero más cercano. Todas las dependencias son de tipo fin-inicio y la holgura se calcula con los métodos de ruta crítica.

{\scriptsize
\setlength{\tabcolsep}{2pt}
\begin{longtable}{|p{0.9cm}|>{\raggedright\arraybackslash}p{4.4cm}|p{0.5cm}|p{0.5cm}|p{0.5cm}|p{0.8cm}|p{0.8cm}|p{0.75cm}|p{1.1cm}|p{1.1cm}|p{1.6cm}|p{0.75cm}|}
\caption{Programación PERT del proyecto}\\
\hline
\textbf{ID} & \textbf{Actividad} & \textbf{O} & \textbf{M} & \textbf{P} & \textbf{E} & $\sigma^2$ & \textbf{Dur.} & \textbf{Inicio} & \textbf{Fin} & \textbf{Predec.} & \textbf{Holg.} \\ \hline
\endfirsthead
\hline
\textbf{ID} & \textbf{Actividad} & \textbf{O} & \textbf{M} & \textbf{P} & \textbf{E} & $\sigma^2$ & \textbf{Dur.} & \textbf{Inicio} & \textbf{Fin} & \textbf{Predec.} & \textbf{Holg.} \\ \hline
\endhead
1.1 & Conformación del equipo, roles, tablero y repositorio & 2 & 3 & 4 & 3.00 & 0.11 & 3 & 24/09 & 26/09 & -- & 0 \\ \hline
1.2 & Especificación de requisitos (ERS) e historias de usuario & 3 & 4 & 6 & 4.17 & 0.25 & 4 & 27/09 & 30/09 & 1.1 & 0 \\ \hline
1.3 & Entornos Docker, CI y credenciales de APIs & 3 & 5 & 7 & 5.00 & 0.44 & 5 & 27/09 & 01/10 & 1.1 & 4 \\ \hline
1.4 & Diseño del esquema de BD (ER) y migraciones & 3 & 5 & 7 & 5.00 & 0.44 & 5 & 01/10 & 05/10 & 1.2 & 0 \\ \hline
1.5 & Estudio de mercado y modelo de negocio (Canvas) & 4 & 6 & 9 & 6.17 & 0.69 & 6 & 27/09 & 02/10 & 1.1 & 39 \\ \hline
2.1 & Autenticación y RBAC & 4 & 6 & 9 & 6.17 & 0.69 & 6 & 06/10 & 11/10 & 1.3, 1.4 & 3 \\ \hline
2.2 & Catálogo (sin gradear, gradeadas y PSA) y publicaciones & 6 & 9 & 14 & 9.33 & 1.78 & 9 & 06/10 & 14/10 & 1.3, 1.4 & 0 \\ \hline
2.3 & Adaptadores de APIs externas y scheduler de precios & 5 & 7 & 10 & 7.17 & 0.69 & 7 & 06/10 & 12/10 & 1.3, 1.4 & 17 \\ \hline
2.4 & Contrato OpenAPI y datos de ejemplo & 2 & 3 & 4 & 3.00 & 0.11 & 3 & 06/10 & 08/10 & 1.4 & 6 \\ \hline
3.1 & Prototipos UI y SPA base (sistema de diseño) & 4 & 6 & 9 & 6.17 & 0.69 & 6 & 09/10 & 14/10 & 2.4, 1.3 & 6 \\ \hline
3.2 & Catálogo por secciones, búsqueda y detalle con gráfica de precios & 5 & 8 & 11 & 8.00 & 1.00 & 8 & 15/10 & 22/10 & 3.1 & 6 \\ \hline
3.3 & Publicación de artículos (4 fotos) y paneles por rol & 5 & 7 & 10 & 7.17 & 0.69 & 7 & 15/10 & 21/10 & 3.1, 2.1 & 8 \\ \hline
4.1 & Mercado Pago (sandbox) y webhooks & 5 & 7 & 10 & 7.17 & 0.69 & 7 & 15/10 & 21/10 & 2.2, 2.1 & 0 \\ \hline
4.2 & Máquina de estados del escrow y disputas & 5 & 7 & 10 & 7.17 & 0.69 & 7 & 22/10 & 28/10 & 4.1 & 0 \\ \hline
4.3 & Geocodificación Mapbox y mapa animado (Leaflet) & 5 & 7 & 11 & 7.33 & 1.00 & 7 & 29/10 & 04/11 & 4.2, 3.2 & 0 \\ \hline
4.4 & Paneles de Recolector y Admin & 4 & 6 & 8 & 6.00 & 0.44 & 6 & 29/10 & 03/11 & 4.2, 3.3 & 1 \\ \hline
5.1 & Pruebas unitarias y de integración & 4 & 6 & 8 & 6.00 & 0.44 & 6 & 29/10 & 03/11 & 4.2, 2.3, 3.2 & 1 \\ \hline
5.2 & Pruebas de seguridad y de carga & 4 & 6 & 9 & 6.17 & 0.69 & 6 & 05/11 & 10/11 & 5.1, 4.3, 4.4 & 0 \\ \hline
5.3 & Corrección de errores y validación con usuarios & 5 & 7 & 10 & 7.17 & 0.69 & 7 & 11/11 & 17/11 & 5.2, 1.5 & 0 \\ \hline
6.1 & Empaquetado en Docker y respaldos & 2 & 3 & 5 & 3.17 & 0.25 & 3 & 11/11 & 13/11 & 5.2 & 4 \\ \hline
6.2 & Documentación técnica y manual de usuario & 3 & 4 & 6 & 4.17 & 0.25 & 4 & 04/11 & 07/11 & 5.1 & 10 \\ \hline
6.3 & Cierre: demostración y entrega final & 2 & 3 & 4 & 3.00 & 0.11 & 3 & 18/11 & 20/11 & 5.3, 6.1, 6.2 & 0 \\ \hline
\end{longtable}
}

\begin{center}
\renewcommand{\arraystretch}{1.3}
\begin{tabular}{@{}p{0.5cm}p{3.6cm}p{0.5cm}p{3.6cm}@{}}
\tikz{\fill[blue!35] (0,0) rectangle (0.4,0.28);} & Trabajo programado de la actividad &
\tikz{\fill[red!60] (0,0) rectangle (0.4,0.28);} & Actividad de la ruta crítica \\
\tikz{\fill[blue!12] (0,0) rectangle (0.4,0.28);} & Holgura: margen en que la actividad podría moverse sin retrasar el proyecto &
\tikz{\fill[gray!45] (0,0) rectangle (0.4,0.28);} & Reserva de contingencia general \\
\tikz{\fill[orange!80] (0.2,0.14) -- (0.4,0.28) -- (0.2,0) -- (0,0.14) -- cycle;} & Hito (entrega final o fecha límite) & & \\
\end{tabular}
\end{center}


\begin{figure}[H]
\centering
\begin{ganttchart}[x unit=1.23cm, y unit title=0.5cm, y unit chart=0.6cm, title height=1, bar height=0.6,
  hgrid, vgrid, title label font=\scriptsize, bar label font=\footnotesize,
  group label font=\footnotesize\bfseries, bar/.append style={fill=blue!35}, group/.append style={fill=black!65},
  milestone/.append style={fill=orange!80}]{1}{7}
\gantttitle{Semana 1: 24/09/2026 -- 30/09/2026}{7} \\
\gantttitlelist{24,25,26,27,28,29,30}{1} \\
\ganttgroup{F1. Inicio y configuración}{1}{7} \\
\ganttbar[bar/.append style={fill=red!60}]{1.1 Equipo, roles y repositorio}{1}{3} \\
\ganttbar[bar/.append style={fill=red!60}]{1.2 ERS e historias de usuario}{4}{7} \\
\ganttbar{1.3 Entornos Docker, CI y APIs}{4}{7} \\
\ganttbar{1.5 Estudio de mercado y Canvas}{4}{7}
\end{ganttchart}
\caption{Semana 1 del cronograma (24/09/2026 -- 30/09/2026)}
\end{figure}

\begin{figure}[H]
\centering
\begin{ganttchart}[x unit=1.23cm, y unit title=0.5cm, y unit chart=0.6cm, title height=1, bar height=0.6,
  hgrid, vgrid, title label font=\scriptsize, bar label font=\footnotesize,
  group label font=\footnotesize\bfseries, bar/.append style={fill=blue!35}, group/.append style={fill=black!65},
  milestone/.append style={fill=orange!80}]{1}{7}
\gantttitle{Semana 2: 01/10/2026 -- 07/10/2026}{7} \\
\gantttitlelist{1,2,3,4,5,6,7}{1} \\
\ganttgroup{F1. Inicio y configuración}{1}{7} \\
\ganttbar{1.3 Entornos Docker, CI y APIs}{1}{1} \\
\ganttbar[bar/.append style={fill=blue!12}]{}{2}{5} \\
\ganttbar[bar/.append style={fill=red!60}]{1.4 Esquema de BD y migraciones}{1}{5} \\
\ganttbar{1.5 Estudio de mercado y Canvas}{1}{2} \\
\ganttbar[bar/.append style={fill=blue!12}]{}{3}{7} \\
\ganttgroup{F2. Backend y APIs}{6}{7} \\
\ganttbar{2.1 Autenticación y RBAC}{6}{7} \\
\ganttbar[bar/.append style={fill=red!60}]{2.2 Catálogo y publicaciones}{6}{7} \\
\ganttbar{2.3 Adaptadores de APIs y precios}{6}{7} \\
\ganttbar{2.4 Contrato OpenAPI y fixtures}{6}{7}
\end{ganttchart}
\caption{Semana 2 del cronograma (01/10/2026 -- 07/10/2026)}
\end{figure}

\begin{figure}[H]
\centering
\begin{ganttchart}[x unit=1.23cm, y unit title=0.5cm, y unit chart=0.6cm, title height=1, bar height=0.6,
  hgrid, vgrid, title label font=\scriptsize, bar label font=\footnotesize,
  group label font=\footnotesize\bfseries, bar/.append style={fill=blue!35}, group/.append style={fill=black!65},
  milestone/.append style={fill=orange!80}]{1}{7}
\gantttitle{Semana 3: 08/10/2026 -- 14/10/2026}{7} \\
\gantttitlelist{8,9,10,11,12,13,14}{1} \\
\ganttgroup{F1. Inicio y configuración}{1}{7} \\
\ganttbar[bar/.append style={fill=blue!12}]{1.5 Estudio de mercado y Canvas}{1}{7} \\
\ganttgroup{F2. Backend y APIs}{1}{7} \\
\ganttbar{2.1 Autenticación y RBAC}{1}{4} \\
\ganttbar[bar/.append style={fill=blue!12}]{}{5}{7} \\
\ganttbar[bar/.append style={fill=red!60}]{2.2 Catálogo y publicaciones}{1}{7} \\
\ganttbar{2.3 Adaptadores de APIs y precios}{1}{5} \\
\ganttbar[bar/.append style={fill=blue!12}]{}{6}{7} \\
\ganttbar{2.4 Contrato OpenAPI y fixtures}{1}{1} \\
\ganttbar[bar/.append style={fill=blue!12}]{}{2}{7} \\
\ganttgroup{F3. Frontend}{2}{7} \\
\ganttbar{3.1 Prototipos UI y SPA base}{2}{7}
\end{ganttchart}
\caption{Semana 3 del cronograma (08/10/2026 -- 14/10/2026)}
\end{figure}

\begin{figure}[H]
\centering
\begin{ganttchart}[x unit=1.23cm, y unit title=0.5cm, y unit chart=0.6cm, title height=1, bar height=0.6,
  hgrid, vgrid, title label font=\scriptsize, bar label font=\footnotesize,
  group label font=\footnotesize\bfseries, bar/.append style={fill=blue!35}, group/.append style={fill=black!65},
  milestone/.append style={fill=orange!80}]{1}{7}
\gantttitle{Semana 4: 15/10/2026 -- 21/10/2026}{7} \\
\gantttitlelist{15,16,17,18,19,20,21}{1} \\
\ganttgroup{F1. Inicio y configuración}{1}{7} \\
\ganttbar[bar/.append style={fill=blue!12}]{1.5 Estudio de mercado y Canvas}{1}{7} \\
\ganttgroup{F2. Backend y APIs}{1}{7} \\
\ganttbar[bar/.append style={fill=blue!12}]{2.3 Adaptadores de APIs y precios}{1}{7} \\
\ganttgroup{F3. Frontend}{1}{7} \\
\ganttbar[bar/.append style={fill=blue!12}]{3.1 Prototipos UI y SPA base}{1}{6} \\
\ganttbar{3.2 Catálogo y gráfica de precios}{1}{7} \\
\ganttbar{3.3 Publicación y paneles por rol}{1}{7} \\
\ganttgroup{F4. Pagos, escrow y mapas}{1}{7} \\
\ganttbar[bar/.append style={fill=red!60}]{4.1 Mercado Pago y webhooks}{1}{7}
\end{ganttchart}
\caption{Semana 4 del cronograma (15/10/2026 -- 21/10/2026)}
\end{figure}

\begin{figure}[H]
\centering
\begin{ganttchart}[x unit=1.23cm, y unit title=0.5cm, y unit chart=0.6cm, title height=1, bar height=0.6,
  hgrid, vgrid, title label font=\scriptsize, bar label font=\footnotesize,
  group label font=\footnotesize\bfseries, bar/.append style={fill=blue!35}, group/.append style={fill=black!65},
  milestone/.append style={fill=orange!80}]{1}{7}
\gantttitle{Semana 5: 22/10/2026 -- 28/10/2026}{7} \\
\gantttitlelist{22,23,24,25,26,27,28}{1} \\
\ganttgroup{F1. Inicio y configuración}{1}{7} \\
\ganttbar[bar/.append style={fill=blue!12}]{1.5 Estudio de mercado y Canvas}{1}{7} \\
\ganttgroup{F2. Backend y APIs}{1}{7} \\
\ganttbar[bar/.append style={fill=blue!12}]{2.3 Adaptadores de APIs y precios}{1}{7} \\
\ganttgroup{F3. Frontend}{1}{7} \\
\ganttbar{3.2 Catálogo y gráfica de precios}{1}{1} \\
\ganttbar[bar/.append style={fill=blue!12}]{}{2}{7} \\
\ganttbar[bar/.append style={fill=blue!12}]{3.3 Publicación y paneles por rol}{1}{7} \\
\ganttgroup{F4. Pagos, escrow y mapas}{1}{7} \\
\ganttbar[bar/.append style={fill=red!60}]{4.2 Escrow y disputas}{1}{7}
\end{ganttchart}
\caption{Semana 5 del cronograma (22/10/2026 -- 28/10/2026)}
\end{figure}

\begin{figure}[H]
\centering
\begin{ganttchart}[x unit=1.23cm, y unit title=0.5cm, y unit chart=0.6cm, title height=1, bar height=0.6,
  hgrid, vgrid, title label font=\scriptsize, bar label font=\footnotesize,
  group label font=\footnotesize\bfseries, bar/.append style={fill=blue!35}, group/.append style={fill=black!65},
  milestone/.append style={fill=orange!80}]{1}{7}
\gantttitle{Semana 6: 29/10/2026 -- 04/11/2026}{7} \\
\gantttitlelist{29,30,31,1,2,3,4}{1} \\
\ganttgroup{F1. Inicio y configuración}{1}{7} \\
\ganttbar[bar/.append style={fill=blue!12}]{1.5 Estudio de mercado y Canvas}{1}{7} \\
\ganttgroup{F2. Backend y APIs}{1}{1} \\
\ganttbar[bar/.append style={fill=blue!12}]{2.3 Adaptadores de APIs y precios}{1}{1} \\
\ganttgroup{F3. Frontend}{1}{1} \\
\ganttbar[bar/.append style={fill=blue!12}]{3.3 Publicación y paneles por rol}{1}{1} \\
\ganttgroup{F4. Pagos, escrow y mapas}{1}{7} \\
\ganttbar[bar/.append style={fill=red!60}]{4.3 Mapbox y mapa animado}{1}{7} \\
\ganttbar{4.4 Paneles Recolector y Admin}{1}{6} \\
\ganttbar[bar/.append style={fill=blue!12}]{}{7}{7} \\
\ganttgroup{F5. Integración y pruebas}{1}{7} \\
\ganttbar{5.1 Pruebas unitarias e integración}{1}{6} \\
\ganttbar[bar/.append style={fill=blue!12}]{}{7}{7} \\
\ganttgroup{F6. Despliegue y cierre}{7}{7} \\
\ganttbar{6.2 Documentación y manual}{7}{7}
\end{ganttchart}
\caption{Semana 6 del cronograma (29/10/2026 -- 04/11/2026)}
\end{figure}

\begin{figure}[H]
\centering
\begin{ganttchart}[x unit=1.23cm, y unit title=0.5cm, y unit chart=0.6cm, title height=1, bar height=0.6,
  hgrid, vgrid, title label font=\scriptsize, bar label font=\footnotesize,
  group label font=\footnotesize\bfseries, bar/.append style={fill=blue!35}, group/.append style={fill=black!65},
  milestone/.append style={fill=orange!80}]{1}{7}
\gantttitle{Semana 7: 05/11/2026 -- 11/11/2026}{7} \\
\gantttitlelist{5,6,7,8,9,10,11}{1} \\
\ganttgroup{F1. Inicio y configuración}{1}{6} \\
\ganttbar[bar/.append style={fill=blue!12}]{1.5 Estudio de mercado y Canvas}{1}{6} \\
\ganttgroup{F5. Integración y pruebas}{1}{7} \\
\ganttbar[bar/.append style={fill=red!60}]{5.2 Pruebas de seguridad y carga}{1}{6} \\
\ganttbar[bar/.append style={fill=red!60}]{5.3 Correcciones y validación}{7}{7} \\
\ganttgroup{F6. Despliegue y cierre}{1}{7} \\
\ganttbar{6.1 Empaquetado Docker y respaldos}{7}{7} \\
\ganttbar{6.2 Documentación y manual}{1}{3} \\
\ganttbar[bar/.append style={fill=blue!12}]{}{4}{7}
\end{ganttchart}
\caption{Semana 7 del cronograma (05/11/2026 -- 11/11/2026)}
\end{figure}

\begin{figure}[H]
\centering
\begin{ganttchart}[x unit=1.23cm, y unit title=0.5cm, y unit chart=0.6cm, title height=1, bar height=0.6,
  hgrid, vgrid, title label font=\scriptsize, bar label font=\footnotesize,
  group label font=\footnotesize\bfseries, bar/.append style={fill=blue!35}, group/.append style={fill=black!65},
  milestone/.append style={fill=orange!80}]{1}{7}
\gantttitle{Semana 8: 12/11/2026 -- 18/11/2026}{7} \\
\gantttitlelist{12,13,14,15,16,17,18}{1} \\
\ganttgroup{F5. Integración y pruebas}{1}{6} \\
\ganttbar[bar/.append style={fill=red!60}]{5.3 Correcciones y validación}{1}{6} \\
\ganttgroup{F6. Despliegue y cierre}{1}{7} \\
\ganttbar{6.1 Empaquetado Docker y respaldos}{1}{2} \\
\ganttbar[bar/.append style={fill=blue!12}]{}{3}{6} \\
\ganttbar[bar/.append style={fill=blue!12}]{6.2 Documentación y manual}{1}{6} \\
\ganttbar[bar/.append style={fill=red!60}]{6.3 Cierre y entrega final}{7}{7}
\end{ganttchart}
\caption{Semana 8 del cronograma (12/11/2026 -- 18/11/2026)}
\end{figure}

\begin{figure}[H]
\centering
\begin{ganttchart}[x unit=1.35cm, y unit title=0.5cm, y unit chart=0.6cm, title height=1, bar height=0.6,
  hgrid, vgrid, title label font=\scriptsize, bar label font=\footnotesize,
  group label font=\footnotesize\bfseries, bar/.append style={fill=blue!35}, group/.append style={fill=black!65},
  milestone/.append style={fill=orange!80}]{1}{6}
\gantttitle{Semana 9: 19/11/2026 -- 24/11/2026}{6} \\
\gantttitlelist{19,20,21,22,23,24}{1} \\
\ganttgroup{F6. Despliegue y cierre}{1}{2} \\
\ganttbar[bar/.append style={fill=red!60}]{6.3 Cierre y entrega final}{1}{2} \\
\ganttbar[bar/.append style={fill=gray!45}]{Reserva de contingencia}{3}{6} \\
\ganttmilestone{Entrega final}{2} \\
\ganttmilestone{Fecha límite}{6}
\end{ganttchart}
\caption{Semana 9 del cronograma (19/11/2026 -- 24/11/2026)}
\end{figure}

\textbf{Ruta crítica}
La ruta crítica está formada por las actividades 1.1 $\rightarrow$ 1.2 $\rightarrow$ 1.4 $\rightarrow$ 2.2 $\rightarrow$ 4.1 $\rightarrow$ 4.2 $\rightarrow$ 4.3 $\rightarrow$ 5.2 $\rightarrow$ 5.3 $\rightarrow$ 6.3. Su duración programada es de 58 días (del 24/09/2026 al 20/11/2026, día 58); su duración esperada sin redondeo es de 59.5 días, con desviación estándar de 2.54 días. Como el plazo es de 62 días, existe una reserva de 4 días (del 21/11/2026 al 24/11/2026) para absorber retrasos, ajustes de alcance o imprevistos técnicos. Con la aproximación normal de PERT, la probabilidad de concluir dentro del plazo es aproximadamente 84\,\%; esta cifra debe tomarse como referencia, porque ignora las rutas paralelas cercanas a la crítica.


\newpage
\section{Plan de trabajo (Base line)}

\label{sec:ci}

\bigskip
\noindent\begin{minipage}{\textwidth}
\small\renewcommand{\arraystretch}{1.9}
\begin{tabularx}{\textwidth}{|p{3.3cm}|Y|}
\hline
Nombre de CI & Anteproyecto (documento) \\ \hline
Identificador & ANTE-1.0 \\ \hline
Descripción de CI & Documento de anteproyecto: objetivos, definición del producto, contexto de negocio, estrategias y plan de trabajo. \\ \hline
Fecha(s) de CI & 21/09/2026 \\ \hline
Tipo de CI & Documentación \\ \hline
Propietario & Líder de proyecto \\ \hline
Relación de CI & Base de MODNEG-1.0 y ERS-1.0. \\ \hline
\end{tabularx}
\end{minipage}\par\bigskip
\noindent\begin{minipage}{\textwidth}
\small\renewcommand{\arraystretch}{1.9}
\begin{tabularx}{\textwidth}{|p{3.3cm}|Y|}
\hline
Nombre de CI & Modelo de negocio y estudio de mercado \\ \hline
Identificador & MODNEG-1.0 \\ \hline
Descripción de CI & Análisis de competencia, encuesta a coleccionistas, Business Model Canvas y modelo de ingresos. \\ \hline
Fecha(s) de CI & 02/10/2026 \\ \hline
Tipo de CI & Documentación \\ \hline
Propietario & Analista de negocio \\ \hline
Relación de CI & Alimenta a ERS-1.0; se contrasta en la actividad 5.3. \\ \hline
\end{tabularx}
\end{minipage}\par\bigskip
\noindent\begin{minipage}{\textwidth}
\small\renewcommand{\arraystretch}{1.9}
\begin{tabularx}{\textwidth}{|p{3.3cm}|Y|}
\hline
Nombre de CI & Especificación de requisitos (ERS) \\ \hline
Identificador & ERS-1.0 \\ \hline
Descripción de CI & Requisitos funcionales y no funcionales, historias de usuario y criterios de aceptación. \\ \hline
Fecha(s) de CI & 30/09/2026 \\ \hline
Tipo de CI & Documentación \\ \hline
Propietario & Analista de negocio \\ \hline
Relación de CI & Derivación: implementado en API-1.0 y WEB-1.0; base de CP-1.0. \\ \hline
\end{tabularx}
\end{minipage}\par\bigskip
\noindent\begin{minipage}{\textwidth}
\small\renewcommand{\arraystretch}{1.9}
\begin{tabularx}{\textwidth}{|p{3.3cm}|Y|}
\hline
Nombre de CI & Repositorio y control de versiones \\ \hline
Identificador & VERSIONCTRL-1.0 \\ \hline
Descripción de CI & Repositorio Git con Git Flow (ramas dev y main), tablero de tareas y protección de ramas. \\ \hline
Fecha(s) de CI & 26/09/2026 \\ \hline
Tipo de CI & Control de versiones \\ \hline
Propietario & Líder de proyecto \\ \hline
Relación de CI & Base de API-1.0, WEB-1.0 y PRECIOS-1.0. \\ \hline
\end{tabularx}
\end{minipage}\par\bigskip
\noindent\begin{minipage}{\textwidth}
\small\renewcommand{\arraystretch}{1.9}
\begin{tabularx}{\textwidth}{|p{3.3cm}|Y|}
\hline
Nombre de CI & Entorno Docker y CI \\ \hline
Identificador & DOCKER-1.0 \\ \hline
Descripción de CI & Docker Compose para levantar el sistema en cualquier equipo del grupo (entorno de desarrollo en Windows) y pipeline de GitHub Actions para pruebas. \\ \hline
Fecha(s) de CI & 01/10/2026 \\ \hline
Tipo de CI & Configuración \\ \hline
Propietario & Arquitecto / Backend \\ \hline
Relación de CI & Depende de VERSIONCTRL-1.0; base de API-1.0 y DEPLOY-1.0. \\ \hline
\end{tabularx}
\end{minipage}\par\bigskip
\noindent\begin{minipage}{\textwidth}
\small\renewcommand{\arraystretch}{1.9}
\begin{tabularx}{\textwidth}{|p{3.3cm}|Y|}
\hline
Nombre de CI & Cuentas y credenciales de APIs \\ \hline
Identificador & APIEXT-1.0 \\ \hline
Descripción de CI & Credenciales y cuentas de sandbox: PokemonPriceTracker, pokemontcg.io, Mercado Libre, Mercado Pago y Mapbox. \\ \hline
Fecha(s) de CI & 01/10/2026 \\ \hline
Tipo de CI & Servicios externos \\ \hline
Propietario & Arquitecto / Backend \\ \hline
Relación de CI & Requerido por PRECIOS-1.0, API-1.0 y WEB-1.0. \\ \hline
\end{tabularx}
\end{minipage}\par\bigskip
\noindent\begin{minipage}{\textwidth}
\small\renewcommand{\arraystretch}{1.9}
\begin{tabularx}{\textwidth}{|p{3.3cm}|Y|}
\hline
Nombre de CI & Esquema de base de datos \\ \hline
Identificador & DBMARIA-1.0 \\ \hline
Descripción de CI & Esquema MariaDB (usuarios, roles, artículos, publicaciones con campos de grado y certificado, órdenes, pagos, envíos, snapshots de precio) y migraciones con Alembic. \\ \hline
Fecha(s) de CI & 05/10/2026 \\ \hline
Tipo de CI & Base de datos \\ \hline
Propietario & Arquitecto / Backend \\ \hline
Relación de CI & Depende de ERS-1.0; base de API-1.0 y OPENAPI-1.0. \\ \hline
\end{tabularx}
\end{minipage}\par\bigskip
\noindent\begin{minipage}{\textwidth}
\small\renewcommand{\arraystretch}{1.9}
\begin{tabularx}{\textwidth}{|p{3.3cm}|Y|}
\hline
Nombre de CI & Contrato de API y datos de ejemplo \\ \hline
Identificador & OPENAPI-1.0 \\ \hline
Descripción de CI & Especificación OpenAPI de los endpoints y fixtures para desarrollar el frontend en paralelo. \\ \hline
Fecha(s) de CI & 08/10/2026 \\ \hline
Tipo de CI & Diseño / Documentación \\ \hline
Propietario & Arquitecto / Backend \\ \hline
Relación de CI & Depende de DBMARIA-1.0; contrato para WEB-1.0 y CP-1.0. \\ \hline
\end{tabularx}
\end{minipage}\par\bigskip
\noindent\begin{minipage}{\textwidth}
\small\renewcommand{\arraystretch}{1.9}
\begin{tabularx}{\textwidth}{|p{3.3cm}|Y|}
\hline
Nombre de CI & Módulo de precios externos \\ \hline
Identificador & PRECIOS-1.0 \\ \hline
Descripción de CI & Adaptadores de APIs externas, scheduler de precios, caché y almacenamiento de snapshots. \\ \hline
Fecha(s) de CI & 12/10/2026 \\ \hline
Tipo de CI & Código fuente \\ \hline
Propietario & Arquitecto / Backend \\ \hline
Relación de CI & Depende de DBMARIA-1.0 y APIEXT-1.0; provee datos a API-1.0. \\ \hline
\end{tabularx}
\end{minipage}\par\bigskip
\noindent\begin{minipage}{\textwidth}
\small\renewcommand{\arraystretch}{1.9}
\begin{tabularx}{\textwidth}{|p{3.3cm}|Y|}
\hline
Nombre de CI & Backend (API REST) \\ \hline
Identificador & API-1.0 \\ \hline
Descripción de CI & FastAPI: autenticación, RBAC, catálogo (sin gradar, gradadas y PSA), publicaciones, pagos en sandbox, escrow, disputas y envíos. \\ \hline
Fecha(s) de CI & 28/10/2026 \\ \hline
Tipo de CI & Código fuente \\ \hline
Propietario & Arquitecto / Backend \\ \hline
Relación de CI & Depende de DBMARIA-1.0, DOCKER-1.0 y APIEXT-1.0; consumido por WEB-1.0. \\ \hline
\end{tabularx}
\end{minipage}\par\bigskip
\noindent\begin{minipage}{\textwidth}
\small\renewcommand{\arraystretch}{1.9}
\begin{tabularx}{\textwidth}{|p{3.3cm}|Y|}
\hline
Nombre de CI & Prototipos de interfaz \\ \hline
Identificador & PROTO-1.0 \\ \hline
Descripción de CI & Wireframes y sistema de diseño (paleta, tipografía y componentes) de la estética retro-moderna. \\ \hline
Fecha(s) de CI & 14/10/2026 \\ \hline
Tipo de CI & Diseño \\ \hline
Propietario & Desarrollador frontend \\ \hline
Relación de CI & Depende de ERS-1.0; guía de WEB-1.0. \\ \hline
\end{tabularx}
\end{minipage}\par\bigskip
\noindent\begin{minipage}{\textwidth}
\small\renewcommand{\arraystretch}{1.9}
\begin{tabularx}{\textwidth}{|p{3.3cm}|Y|}
\hline
Nombre de CI & Frontend (SPA) \\ \hline
Identificador & WEB-1.0 \\ \hline
Descripción de CI & SPA React: catálogo por secciones (con sección PSA), detalle con gráficas, publicación, paneles por rol y mapa animado. \\ \hline
Fecha(s) de CI & 03/11/2026 \\ \hline
Tipo de CI & Código fuente \\ \hline
Propietario & Desarrollador frontend \\ \hline
Relación de CI & Depende de PROTO-1.0, OPENAPI-1.0 y API-1.0. \\ \hline
\end{tabularx}
\end{minipage}\par\bigskip
\noindent\begin{minipage}{\textwidth}
\small\renewcommand{\arraystretch}{1.9}
\begin{tabularx}{\textwidth}{|p{3.3cm}|Y|}
\hline
Nombre de CI & Casos de prueba \\ \hline
Identificador & CP-1.0 \\ \hline
Descripción de CI & Pruebas unitarias y de integración (pytest) de permisos por rol, escrow y adaptadores de APIs. \\ \hline
Fecha(s) de CI & 03/11/2026 \\ \hline
Tipo de CI & Pruebas \\ \hline
Propietario & QA (Tester) \\ \hline
Relación de CI & Valida API-1.0, PRECIOS-1.0 y WEB-1.0. \\ \hline
\end{tabularx}
\end{minipage}\par\bigskip
\noindent\begin{minipage}{\textwidth}
\small\renewcommand{\arraystretch}{1.9}
\begin{tabularx}{\textwidth}{|p{3.3cm}|Y|}
\hline
Nombre de CI & Reporte de seguridad y carga \\ \hline
Identificador & SEG-1.0 \\ \hline
Descripción de CI & Resultados de pruebas de carga (Locust o k6) y revisión de seguridad con base en OWASP Top 10. \\ \hline
Fecha(s) de CI & 10/11/2026 \\ \hline
Tipo de CI & Pruebas \\ \hline
Propietario & QA (Tester) \\ \hline
Relación de CI & Depende de CP-1.0; valida API-1.0 y DEPLOY-1.0. \\ \hline
\end{tabularx}
\end{minipage}\par\bigskip
\noindent\begin{minipage}{\textwidth}
\small\renewcommand{\arraystretch}{1.9}
\begin{tabularx}{\textwidth}{|p{3.3cm}|Y|}
\hline
Nombre de CI & Despliegue y respaldos \\ \hline
Identificador & DEPLOY-1.0 \\ \hline
Descripción de CI & Empaquetado del sistema con Docker Compose en un entorno controlado (equipo de desarrollo con Windows) y respaldos automáticos de base de datos e imágenes. \\ \hline
Fecha(s) de CI & 13/11/2026 \\ \hline
Tipo de CI & Configuración \\ \hline
Propietario & Arquitecto / Backend \\ \hline
Relación de CI & Depende de DOCKER-1.0 y SEG-1.0. \\ \hline
\end{tabularx}
\end{minipage}\par\bigskip
\noindent\begin{minipage}{\textwidth}
\small\renewcommand{\arraystretch}{1.9}
\begin{tabularx}{\textwidth}{|p{3.3cm}|Y|}
\hline
Nombre de CI & Manual y documentación técnica \\ \hline
Identificador & MANUAL-1.0 \\ \hline
Descripción de CI & Manual de usuario por rol, guía de despliegue y documentación técnica. \\ \hline
Fecha(s) de CI & 07/11/2026 \\ \hline
Tipo de CI & Documentación \\ \hline
Propietario & Analista de negocio \\ \hline
Relación de CI & Describe WEB-1.0, API-1.0 y DEPLOY-1.0. \\ \hline
\end{tabularx}
\end{minipage}\par\bigskip
\noindent\begin{minipage}{\textwidth}
\small\renewcommand{\arraystretch}{1.9}
\begin{tabularx}{\textwidth}{|p{3.3cm}|Y|}
\hline
Nombre de CI & Entrega final \\ \hline
Identificador & ENTREGA-1.0 \\ \hline
Descripción de CI & Presentación y demostración final con todos los elementos de configuración integrados. \\ \hline
Fecha(s) de CI & 20/11/2026 \\ \hline
Tipo de CI & Entrega \\ \hline
Propietario & Líder de proyecto \\ \hline
Relación de CI & Integra todos los CI anteriores. \\ \hline
\end{tabularx}
\end{minipage}\par\bigskip


\subsection{Planificaciones}

\subsection{Cronograma de trabajo con hitos clave}

\begin{longtable}{|>{\raggedright\arraybackslash}p{3.1cm}|>{\raggedright\arraybackslash}p{7.1cm}|>{\raggedright\arraybackslash}p{3.7cm}|}
\caption{Cronograma por fases e hitos}\\
\hline
\textbf{Fase} & \textbf{Hito clave} & \textbf{Fecha estimada} \\ \hline
\endfirsthead
\hline
\textbf{Fase} & \textbf{Hito clave} & \textbf{Fecha estimada} \\ \hline
\endhead
\textbf{1. Inicio y configuración}\newline (24/09 -- 05/10/2026) & Conformación del equipo, roles, tablero y repositorio (\texttt{VERSIONCTRL-1.0}) & \mbox{24/09 -- 26/09/2026} \\
\cline{2-3}
 & Especificación de requisitos (ERS) e historias de usuario (\texttt{ERS-1.0}) & \mbox{27/09 -- 30/09/2026} \\
\cline{2-3}
 & Entornos Docker, CI y credenciales de APIs (\texttt{DOCKER-1.0, APIEXT-1.0}) & \mbox{27/09 -- 01/10/2026} \\
\cline{2-3}
 & Diseño del esquema de BD (ER) y migraciones (\texttt{DBMARIA-1.0}) & \mbox{01/10 -- 05/10/2026} \\
\cline{2-3}
 & Estudio de mercado y modelo de negocio (Canvas) (\texttt{MODNEG-1.0}) & \mbox{27/09 -- 02/10/2026} \\
\hline
\textbf{2. Backend y APIs}\newline (06/10 -- 14/10/2026) & Autenticación y RBAC & \mbox{06/10 -- 11/10/2026} \\
\cline{2-3}
 & Catálogo (sin gradar, gradadas y PSA) y publicaciones (\texttt{API-1.0 (parcial)}) & \mbox{06/10 -- 14/10/2026} \\
\cline{2-3}
 & Adaptadores de APIs externas y scheduler de precios (\texttt{PRECIOS-1.0}) & \mbox{06/10 -- 12/10/2026} \\
\cline{2-3}
 & Contrato OpenAPI y datos de ejemplo (\texttt{OPENAPI-1.0}) & \mbox{06/10 -- 08/10/2026} \\
\hline
\textbf{3. Frontend}\newline (09/10 -- 22/10/2026) & Prototipos UI y SPA base (sistema de diseño) (\texttt{PROTO-1.0}) & \mbox{09/10 -- 14/10/2026} \\
\cline{2-3}
 & Catálogo por secciones, búsqueda y detalle con gráfica de precios & \mbox{15/10 -- 22/10/2026} \\
\cline{2-3}
 & Publicación de artículos (4 fotos) y paneles por rol & \mbox{15/10 -- 21/10/2026} \\
\hline
\textbf{4. Pagos, escrow y mapas}\newline (15/10 -- 04/11/2026) & Mercado Pago (sandbox) y webhooks & \mbox{15/10 -- 21/10/2026} \\
\cline{2-3}
 & Máquina de estados del escrow y disputas (\texttt{API-1.0}) & \mbox{22/10 -- 28/10/2026} \\
\cline{2-3}
 & Geocodificación Mapbox y mapa animado (Leaflet) & \mbox{29/10 -- 04/11/2026} \\
\cline{2-3}
 & Paneles de Recolector y Admin (\texttt{WEB-1.0}) & \mbox{29/10 -- 03/11/2026} \\
\hline
\textbf{5. Integración y pruebas}\newline (29/10 -- 17/11/2026) & Pruebas unitarias y de integración (\texttt{CP-1.0}) & \mbox{29/10 -- 03/11/2026} \\
\cline{2-3}
 & Pruebas de seguridad y de carga (\texttt{SEG-1.0}) & \mbox{05/11 -- 10/11/2026} \\
\cline{2-3}
 & Corrección de errores y validación con usuarios & \mbox{11/11 -- 17/11/2026} \\
\hline
\textbf{6. Despliegue y cierre}\newline (04/11 -- 20/11/2026) & Empaquetado en Docker y respaldos (\texttt{DEPLOY-1.0}) & \mbox{11/11 -- 13/11/2026} \\
\cline{2-3}
 & Documentación técnica y manual de usuario (\texttt{MANUAL-1.0}) & \mbox{04/11 -- 07/11/2026} \\
\cline{2-3}
 & Cierre: demostración y entrega final (\texttt{ENTREGA-1.0}) & \mbox{18/11 -- 20/11/2026} \\
\hline
\end{longtable}


\newpage


% Impresión Referencias
\nocite{*}
\printbibliography[heading=bibintoc, title={Referencias}]

\end{document}